Полученные результаты позволяют читать модель не как формулу «число агентов
делит срок», а как модель планирования ограниченной человеко-агентной ячейки.

\subsection{Bound, Optimum, and a Concrete Schedule}

$B_4$ отвечает на вопрос, быстрее какого срока закончить невозможно из-за
суммарной загрузки потоков, исходного критического пути или единственного
человеческого ресурса. $T^*$ отвечает на другой вопрос: каков лучший срок с
учётом всех ресурсных порядков. Наконец, $T(\pi)$ характеризует выбранную
политику или фактическое расписание. Равенство $T^*=B_4$ требует сертификата;
в противном случае одинаковая граница не означает одинаковый срок.

\subsection{Resource Chains as a Gap Mechanism}

Повторные обращения нескольких задач к одному разработчику связывают фазы,
которые не лежат на одном пути исходного DAG. Удержание агентного потока во
время ожидания переносит эту конкуренцию дальше по расписанию. Ресурсно-
дополненный фазовый граф делает механизм видимым: его максимальный путь
содержит не только технологические, но и ресурсные дуги. Это объясняет, почему
одни лишь $A,H,W_4,L_4,B_4$ не идентифицируют $T^*$.

\subsection{Parallelism as a Configuration Choice}

Следует различать доступную и настроенную мощность. При неизменных длительностях
дополнительный поток можно оставить свободным, поэтому оптимум не ухудшается.
Но организационный переход к большему $P$ может повысить стоимость интеграции
$C(P)$ и переключений $\gamma(P)$. Тогда сравниваются уже разные процессы, и
конечный минимум возможен. Практически это означает, что $P$ нельзя выбирать в
отрыве от назначения режимов и функций издержек. Формальная конфигурация
$\sigma=(P,\rho)$ фиксирует число потоков и режимы задач; декомпозиция, DAG,
фазовые профили и критерий приёмки образуют остальные условия полностью
заданного сценария. Если вместе с $P$ меняется хотя бы одно из них, наблюдаемая
разность относится к сравнению сценариев целиком, а не к изолированному эффекту
дополнительного потока.

\subsection{Decomposition and Parallelism}

Декомпозиция полезна, пока она сокращает длинные последовательные участки и
создаёт готовую работу для свободных потоков. После этого дополнительные
границы увеличивают спецификацию, интеграцию и число вмешательств человека.
В исследованном операторе декомпозиции поэтому оптимальная гранулярность зависит
от $P$: без доступного параллелизма мелкие независимые задачи почти не дают
выигрыша, а при большом $P$ именно гранулярность может ограничивать
использование мощности. Эксперименты поддерживают такую взаимодополняемость на
выбранных сетках, но не доказывают её для произвольного DAG или правила
дробления; её направление и величина требуют калибровки на реальных процессах.

\subsection{Volume and Placement of Human Work}

Общий $H$ задаёт безусловную последовательную нагрузку и потолок ускорения.
Его расположение --- отдельная характеристика. При одинаковом масштабировании
агентных и человеческих фаз перенос $H$ между задачами не обязан менять
агрегатную границу; при $\gamma>C$ человеческая работа на критическом пути
дороже агентной и удлиняет этот путь. Даже при неизменном $L_4$ фазовый порядок
может менять очередь и $T^*$. Поэтому для проектирования процесса недостаточно
сократить часы ревью: важно также решить, когда и в каких задачах они нужны.

\subsection{Relation to Prior Work}

Потолок по человеческому ресурсу является структурным переносом логики закона
Амдала: последовательной долей здесь служит не участок вычислительной
программы, а измеренное участие разработчика \cite{amdahl1967validity}. Внутренний
минимум по настроенному $P$ согласуется с универсальным законом масштабируемости
в том ограниченном смысле, что растущие конкуренция и согласование способны
сменить насыщение отрицательной отдачей \cite{gunther2008general}. Численная
форма $C(P)$ для агентов при этом остаётся гипотезой модели, а не следствием
этих классических результатов.

Разрыв между $B_4$ и $T^*$ связывает постановку с теорией расписаний и RCPSP:
исходного графа зависимостей задач и суммарной нагрузки недостаточно, когда единственный
обслуживающий ресурс создаёт дополнительные ресурсные цепочки
\cite{graham1966bounds,brucker1999resource}. Дуги ресурсного порядка и общий
сервер являются стандартными механизмами теории расписаний
\cite{pinedo2008scheduling,hall2000parallel}. Модели загрузки и выгрузки уже
содержат сервис до и после обработки и могут удерживать машину до выгрузки
\cite{elidrissi2023loading}; возвратная модель направляет работу после удалённого
сервера на ту же основную машину, но освобождает её на время сервиса
\cite{chakhlevitch2009reentrant}. Следовательно, новизна не заявляется ни для
DAG, ни для общего ресурса, повторного сервиса или ресурсно-дополненного графа.

Близкая идея конечного числа эффективно контролируемых исполнителей возникает в
моделях числа роботов на одного оператора, где интервалы обслуживания должны укладываться в окна
автономной работы остальных систем \cite{olsen2004fanout,crandall2005validating}.
Эксперимент Cummings и Mitchell дополнительно показывает, что очередь,
ожидание восстановления ситуационной осведомлённости и переориентация снижают расчётную пропускную
способность оператора, но численно переносить результат имитации UAV на
разработку нельзя \cite{cummings2008predicting}. Среди агентов программирования
MAGIS реализует автоматизированный цикл «разработка --- контроль качества ---
доработка», а CAID --- конкурентные рабочие деревья с планом зависимостей,
слиянием и тестированием \cite{tao2024magis,geng2026caid}. Их показатели
качества и успешности не свидетельствуют о сокращении makespan; в основных
сравнениях CAID календарное время выполнения и стоимость росли вместе с
показателем качества.

Настоящая модель объединяет известные конструкции в отображении на программный
процесс: коэффициент $k$, зависящий от задачи и режима, раздельные $a_i$ и
$h_i$, повторные человеческие фазы, локальное удержание агентного потока,
интеграционные издержки и издержки переключения, а также makespan принятого
результата. В основном корпусе, сформированном по поисковому протоколу с датой
завершения 2026-08-09, не обнаружена их совместная эмпирическая калибровка на
программных задачах; это не универсальное утверждение обо всей литературе. Сама
статья также не закрывает этот эмпирический пробел: вычислительные результаты
уточняют механизм процессного ограничения на синтетических сценариях, но не
заменяют измерения конкретных инструментов и команд.

\subsection{Using the Model}

Сначала следует оценить задачи и режимы в общей шкале, построить DAG и фазовые
профили, затем вычислить $B_4$ как быструю диагностику узкого места. Для
нескольких конкурирующих конфигураций строится точное или эвристическое
расписание и явно фиксируется, получен ли $T^*$ или только $T(\pi)$. Наконец,
проверяется чувствительность к $k$, $h$, $C$ и $\gamma$. Модель полезнее для
сравнения вариантов организации работы, чем как обещание точного календарного
срока без локальной калибровки. Активная ветвь $B_4$ характеризует ограничение
нижней границы, но сама по себе не доказывает, что фактический $T^*$ определяется
той же ветвью: для такого вывода требуется достижимость границы либо анализ
ресурсно-дополненного расписания. Аналогично, автономная агентная фаза означает
только отсутствие участия разработчика на данном интервале; задача с такой
фазой остаётся делегированной, если постановка или приёмка требуют человека.
